\documentclass[12pt, a4paper]{article}

% --- Подключение пакетов ---
\usepackage[utf8]{inputenc}
\usepackage[T2A]{fontenc}
\usepackage[russian]{babel}
\usepackage{amsmath, amsfonts, amssymb} % Математические формулы
\usepackage{graphicx} % Для вставки картинок
\usepackage{geometry} % Настройка полей
\geometry{top=2cm, bottom=2cm, left=3cm, right=1.5cm}
\usepackage{hyperref} % Гиперссылки
\usepackage{booktabs} % Красивые таблицы
\usepackage{array}    % Улучшенные таблицы
\usepackage[style=numeric, backend=biber]{biblatex} % Библиография

% Подключение файла с литературой
\addbibresource{references.bib}

% --- Заголовок ---
\title{\textbf{Литературный обзор: Современные методы прайсинга и хеджирования деривативов с помощью глубокого обучения}}
\author{Бирюков Никита}
\date{\today}

\begin{document}

\maketitle
\tableofcontents
\newpage

% =================================================================
% РАЗДЕЛ 1: ВВЕДЕНИЕ И ОБЗОР ИССЛЕДОВАТЕЛЬСКИХ ГРУПП
% =================================================================
\section{Введение: Смена парадигмы в количественных финансах}

Классическая теория оценки производных финансовых инструментов (деривативов), фундаментально обоснованная работами Блэка, Шоулза и Мертона (1973) \cite{BlackScholes1973}, опирается на стохастическое исчисление и решение уравнений в частных производных (PDE). Традиционные численные методы, такие как метод Монте-Карло (Monte Carlo) и конечно-разностные методы (Finite Difference Methods, FDM), являются «золотым стандартом» для задач низкой размерности. Однако современная финансовая индустрия сталкивается с вызовами, которые делают классические подходы вычислительно неэффективными:
\begin{enumerate}
    \item \textbf{«Проклятие размерности» (Curse of Dimensionality):} Сложность сеточных методов растет экспоненциально ($O(n^d)$). Оценка опциона на корзину из 50--100 акций или расчет CVA (Credit Valuation Adjustment) для портфеля банка классическими методами требует неприемлемых временных затрат.
    \item \textbf{Сложность калибровки (Calibration Bottleneck):} Современные модели, такие как модели грубой волатильности (Rough Volatility), лучше описывают рынок, но не имеют аналитических решений. Их калибровка через Монте-Карло занимает часы, что делает их неприменимыми для реал-тайм трейдинга.
    \item \textbf{Рыночные трения:} Классические модели предполагают идеальный рынок. Учет транзакционных издержек и влияния на рынок (market impact) требует перехода к задачам стохастического управления, которые крайне сложны для аналитического решения.
\end{enumerate}

В ответ на эти вызовы сформировалось направление \textit{Scientific Machine Learning}, предлагающее использование глубоких нейронных сетей (DNN) как универсальных аппроксиматоров для решения финансовых задач.

\subsection{Ведущие исследовательские группы и их вклад}

Анализ современной литературы позволяет выделить несколько ключевых центров компетенций, каждый из которых специализируется на своем классе задач.

% --- Группа 1: Deep BSDE ---
\subsubsection{Группа Deep BSDE (Принстон, ETH Zurich)}
\textbf{Ключевые исследователи:} Weinan E, Jiequn Han, Arnulf Jentzen.

\begin{itemize}
    \item \textbf{Направление исследований:} Решение высокоразмерных параболических PDE (уравнение теплопроводности, Блэка-Шоулза, Аллена-Кана) с помощью нейронных сетей.
    \item \textbf{Достижения:}
    \begin{itemize}
        \item Разработали метод \textit{Deep BSDE} \cite{Han2018}, основанный на стохастической формулировке уравнений (Backward Stochastic Differential Equations).
        \item Впервые успешно решили нелинейное уравнение Блэка-Шоулза с риском дефолта для размерности $d=100$ за разумное время.
        \item Доказали (теоретически и эмпирически), что нейросети могут преодолевать проклятие размерности, аппроксимируя градиент решения без построения пространственной сетки.
    \end{itemize}
    \item \textbf{Оборудование:} В оригинальной статье использовались кластеры с GPU (NVIDIA Tesla K80/P100). Обучение модели для $d=100$ занимало порядка 30--40 минут.
    \item \textbf{Open Source / Код:}
    \begin{itemize}
        \item Официальный репозиторий: \url{https://github.com/frankhan91/DeepBSDE} (TensorFlow).
        \item Существует множество современных портов на PyTorch.
    \end{itemize}
    \item \textbf{Ключевые публикации:} \cite{Han2018}, \cite{Beck2019}.
\end{itemize}

% --- Группа 2: DGM ---
\subsubsection{Группа DGM (University of Oxford, UIUC)}
\textbf{Ключевые исследователи:} Justin Sirignano, Konstantinos Spiliopoulos.

\begin{itemize}
    \item \textbf{Направление исследований:} Бессеточные методы (Meshfree methods) для решения PDE со свободной границей (Free Boundary Problems), что критически важно для оценки Американских опционов.
    \item \textbf{Достижения:}
    \begin{itemize}
        \item Разработали метод \textit{DGM (Deep Galerkin Method)} \cite{Sirignano2018}, который минимизирует невязку дифференциального оператора в случайно выбранных точках пространства.
        \item Успешно решили задачу прайсинга Американских опционов в размерности до $d=200$, что невозможно для традиционных методов.
        \item Предложили метод Монте-Карло для оценки вторых производных (Гессиана), снижающий вычислительную сложность с $O(d^2)$ до $O(d)$.
    \end{itemize}
    \item \textbf{Оборудование:} Использовали суперкомпьютер \textit{Blue Waters} (NCSA) с тысячами узлов GPU, так как метод DGM требует значительных вычислительных ресурсов для сходимости в высоких размерностях.
    \item \textbf{Open Source / Код:}
    \begin{itemize}
        \item Реализации доступны в библиотеке \textit{DeepXDE}: \url{https://github.com/lululxvi/deepxde}.
    \end{itemize}
    \item \textbf{Ключевые публикации:} \cite{Sirignano2018}.
\end{itemize}

% --- Группа 3: Deep Hedging ---
\subsubsection{Группа Deep Hedging (JP Morgan, ETH Zurich)}
\textbf{Ключевые исследователи:} Hans Buehler, Lukas Gonon, Josef Teichmann, Ben Wood.

\begin{itemize}
    \item \textbf{Направление исследований:} Прямая оптимизация торговых стратегий хеджирования в условиях рыночных трений (транзакционные издержки, ликвидность).
    \item \textbf{Достижения:}
    \begin{itemize}
        \item Сформулировали задачу хеджирования не как поиск «справедливой цены», а как минимизацию выпуклой меры риска (Convex Risk Measure) \cite{Buehler2019}.
        \item Показали, что нейросети (LSTM/RNN) способны самостоятельно выучивать сложные стратегии (например, «wait-and-see» или бандинг), которые экономически эффективнее классической Дельта-стратегии при наличии комиссий.
        \item Внедрили данную технологию в реальные бизнес-процессы банка JP Morgan.
    \end{itemize}
    \item \textbf{Оборудование:} Для экспериментов использовались стандартные GPU (например, NVIDIA Tesla V100). Обучение на исторических данных занимает часы, но инференс происходит мгновенно.
    \item \textbf{Open Source / Код:}
    \begin{itemize}
        \item Библиотека PFNet: \url{https://github.com/pfnet-research/deep-hedging}.
    \end{itemize}
    \item \textbf{Ключевые публикации:} \cite{Buehler2019}, \cite{Cuchiero2021} (теоретическое обоснование).
\end{itemize}

% --- Группа 4: Deep Calibration ---
\subsubsection{Группа Deep Calibration (King's College London, Imperial College)}
\textbf{Ключевые исследователи:} Blanka Horvath, Christian Bayer, Aitor Muguruza.

\begin{itemize}
    \item \textbf{Направление исследований:} Ускорение калибровки сложных стохастических моделей, в частности моделей грубой волатильности (Rough Volatility models: rBergomi, rHeston).
    \item \textbf{Достижения:}
    \begin{itemize}
        \item Предложили двухэтапный подход: обучение нейросети-суррогата на синтетических данных (offline) и использование её для мгновенной калибровки (online) \cite{Horvath2019, Bayer2019}.
        \item Достигли ускорения процесса калибровки в $10\,000$ -- $60\,000$ раз по сравнению с традиционными методами, сохраняя точность на уровне метода Монте-Карло.
    \end{itemize}
    \item \textbf{Оборудование:} Обучение проводилось на стандартных GPU (Google Colab, NVIDIA GTX 1080). Генерация обучающей выборки является самой ресурсоемкой частью (требует CPU-кластера или времени), само обучение сети происходит быстро.
    \item \textbf{Open Source / Код:}
    \begin{itemize}
        \item Репозиторий авторов: \url{https://github.com/amuguruza/NN-StochVol-Calibrations}. Содержит готовые ноутбуки для моделей Heston и Rough Bergomi.
    \end{itemize}
    \item \textbf{Ключевые публикации:} \cite{Horvath2019}, \cite{Bayer2019}.
\end{itemize}

% --- Группа 5: Инженерные аспекты и CaNN ---
\subsubsection{Группа прикладных исследований (TU Delft, CWI, NYU)}
\textbf{Ключевые исследователи:} Cornelis Oosterlee, Shuaiqiang Liu (CaNN); Andrey Itkin.

\begin{itemize}
    \item \textbf{Направление исследований:} Инженерная оптимизация нейросетевых солверов, обеспечение отсутствия арбитража, альтернативные методы калибровки.
    \item \textbf{Достижения:}
    \begin{itemize}
        \item Разработали фреймворк \textit{CaNN (Calibration Neural Network)} \cite{Liu2019}, где калибровка параметров модели выполняется через обратное распространение ошибки (Backpropagation) на входной слой сети, без внешнего оптимизатора.
        \item Andrey Itkin предложил методы регуляризации (Soft Constraints) для устранения статического арбитража в предсказаниях нейросети и обосновал необходимость использования гладких функций активации (Softplus, ELU) для корректного расчета «Греков» \cite{Itkin2019}.
    \end{itemize}
    \item \textbf{Open Source / Код:}
    \begin{itemize}
        \item Примеры реализации CaNN доступны в сообществе (поиск по ключевым словам \textit{Neural Network Calibration}).
    \end{itemize}
    \item \textbf{Ключевые публикации:} \cite{Liu2019}, \cite{Itkin2019}.
\end{itemize}

% =================================================================
% РАЗДЕЛ 2: ТЕХНОЛОГИЧЕСКИЙ СТЕК И АРХИТЕКТУРЫ
% =================================================================
\section{Архитектуры нейронных сетей в финансовом моделировании}

Выбор архитектуры нейронной сети является критическим этапом в задачах количественных финансов. В отличие от задач компьютерного зрения или NLP, где существуют устоявшиеся стандарты (CNN, Transformers), в финансовом моделировании архитектура диктуется математической структурой задачи: марковостью процесса, зависимостью от траектории (path-dependency) и требованиями к гладкости производных.

Ниже приведен детальный анализ трех основных классов архитектур, применяемых в прайсинге и хеджировании.

% --- 2.1 FNN ---
\subsection{Полносвязные сети (Feed-Forward Networks / MLP)}

\subsubsection{Теоретическая справка}
Полносвязная сеть (Multilayer Perceptron, MLP) — это базовая архитектура, где каждый нейрон слоя $l$ соединен со всеми нейронами слоя $l+1$. Математически выход слоя описывается как:
\begin{equation}
    h^{(l+1)} = \sigma \left( W^{(l)} h^{(l)} + b^{(l)} \right),
\end{equation}
где $W$ — матрица весов, $b$ — вектор смещения, $\sigma$ — нелинейная функция активации.
Теоретическим обоснованием использования MLP служит \textit{Универсальная теорема аппроксимации} (Hornik et al., 1989), утверждающая, что MLP с одним скрытым слоем может аппроксимировать любую непрерывную функцию с заданной точностью.

\subsubsection{Применение в финансах}
\begin{itemize}
    \item \textbf{Deep Calibration:} Используется как суррогатная модель для отображения «Параметры модели $\to$ Цена опциона» или «Параметры $\to$ Поверхность волатильности» \cite{Horvath2019}.
    \item \textbf{Inverse Problems:} Решение обратных задач, где по цене нужно восстановить параметры (например, локальную волатильность).
\end{itemize}

\subsubsection{Преимущества и недостатки}
\textbf{Плюсы:} Простота реализации; высокая скорость инференса; хорошо изученные методы регуляризации. \\
\textbf{Минусы:} Плохо работают с временными рядами и зависимостями от траектории; подвержены затуханию градиента при большом количестве слоев.

% --- 2.2 ResNet ---
\subsection{Остаточные сети (Residual Networks, ResNets)}

\subsubsection{Теоретическая справка}
ResNets решают проблему затухания градиента в глубоких сетях путем введения «скип-соединений» (skip connections). Вместо того чтобы учить полное отображение $H(x)$, сеть учит остаточную функцию $F(x) = H(x) - x$. Слой ResNet выглядит так:
\begin{equation}
    x_{l+1} = \sigma \left( W_l x_l + b_l \right) + x_l
\end{equation}
Это позволяет сигналу беспрепятственно проходить через сотни слоев, что критически важно для аппроксимации сложных функционалов.

\subsubsection{Применение в финансах}
\begin{itemize}
    \item \textbf{Решение PDE (DGM, Deep BSDE):} Методы DGM \cite{Sirignano2018} и Deep BSDE \cite{Han2018} требуют очень глубоких сетей (до 50 слоев) для точной аппроксимации градиента решения в высокой размерности. Без ResNet-блоков такие сети практически невозможно обучить.
\end{itemize}

\subsubsection{Преимущества и недостатки}
\textbf{Плюсы:} Позволяют строить очень глубокие модели; эффективно обучаются; обеспечивают высокую точность аппроксимации сложных поверхностей. \\
\textbf{Минусы:} Сложнее в настройке гиперпараметров; требуют больше памяти при обучении по сравнению с простым MLP.

% --- 2.3 RNN/LSTM ---
\subsection{Рекуррентные сети (RNN / LSTM)}

\subsubsection{Теоретическая справка}
Рекуррентные сети предназначены для обработки последовательностей. В отличие от FNN, они имеют «внутреннюю память» (hidden state $h_t$), которая передается от шага к шагу:
\begin{equation}
    h_t = \sigma(W_x x_t + W_h h_{t-1} + b).
\end{equation}
Архитектура \textbf{LSTM} (Long Short-Term Memory) добавляет механизм «ворот» (gates), позволяющий сети учить долгосрочные зависимости и избегать взрыва градиента.

\subsubsection{Применение в финансах}
\begin{itemize}
    \item \textbf{Deep Hedging:} В задачах с транзакционными издержками оптимальное действие сегодня зависит от портфеля вчера. Это классическая марковская зависимость, идеально моделируемая LSTM \cite{Buehler2019}.
    \item \textbf{Path-dependent опционы:} Оценка Азиатских опционов или Lookback-опционов, где выплата зависит от всей траектории цены.
    \item \textbf{Прогнозирование временных рядов:} Предсказание реализованной волатильности.
\end{itemize}

\subsubsection{Преимущества и недостатки}
\textbf{Плюсы:} Естественным образом работают с временными рядами; учитывают историю процесса (memory effect); необходимы для моделей с трениями. \\
\textbf{Минусы:} Медленное обучение (плохо распараллеливаются во времени); сложны в интерпретации; требуют больших вычислительных ресурсов (VRAM).

% --- 2.4 СРАВНЕНИЕ ---
\subsection{Сравнительный анализ архитектур}

В Таблице \ref{tab:arch_comparison} представлено сравнение архитектур в контексте задач дипломной работы.

\begin{table}[h!]
    \centering
    \caption{Сравнение архитектур нейронных сетей для финансового моделирования}
    \label{tab:arch_comparison}
    \begin{tabular}{|p{3cm}|p{4cm}|p{4cm}|p{3.5cm}|}
        \hline
        \textbf{Архитектура} & \textbf{Ключевая особенность} & \textbf{Область применения} & \textbf{Сложность обучения} \\
        \hline
        \textbf{FNN (MLP)} & Универсальная аппроксимация, простота & Deep Calibration, простые суррогаты & Низкая (CPU/GPU) \\
        \hline
        \textbf{ResNet} & Skip connections, большая глубина & Решение PDE (DGM, Deep BSDE), высокая размерность & Средняя/Высокая (GPU) \\
        \hline
        \textbf{LSTM/GRU} & Память, обработка последовательностей & Deep Hedging, Path-dependent опционы & Высокая (требует мощных GPU) \\
        \hline
    \end{tabular}
\end{table}

\textbf{Вывод:} Для задач калибровки (Deep Calibration), где требуется мгновенное отображение параметров в цену, оптимальным выбором является FNN. Однако для задач хеджирования (Deep Hedging) и решения PDE в высокой размерности необходимо использование более сложных архитектур типа LSTM и ResNet соответственно, так как они способны улавливать временные зависимости и сложную геометрию решений.

% =================================================================
% ДЕТАЛЬНЫЙ АНАЛИЗ СТАТЕЙ
% =================================================================
\section{Детальный анализ ключевых исследований}

В данном разделе проводится глубокий анализ фундаментальных работ, формирующих методологическую базу исследования.

% --- СТАТЬЯ 1: HAN ET AL. (DEEP BSDE) ---
\subsection{Метод Deep BSDE: Стохастическая оптимизация для высокоразмерных PDE}
\textbf{Статья:} \textit{Solving high-dimensional partial differential equations using deep learning} \cite{Han2018}.

\subsubsection{Проблема и мотивация}
Классические сеточные методы (Finite Difference Methods) страдают от «проклятия размерности»: количество узлов сетки растет как $N^d$, где $d$ — количество активов. Это делает невозможным их применение для $d > 5$. С другой стороны, стандартный метод Монте-Карло, хотя и работает в высоких размерностях, применим только для \textit{линейных} PDE (через формулу Фейнмана-Каца).

Однако многие современные финансовые задачи (например, оценка CVA, прайсинг с учетом дефолтного риска или различных ставок заимствования) описываются \textbf{нелинейными} PDE. Метод Deep BSDE, предложенный группой исследователей из Принстона и ETH Zurich (Weinan E, J. Han, A. Jentzen), стал первым алгоритмом, позволившим решать такие уравнения в размерности $d=100$ за полиномиальное время.

\subsubsection{Математическая теория}
Метод базируется на теории обратных стохастических дифференциальных уравнений (Backward Stochastic Differential Equations, BSDE). Рассмотрим полулинейное параболическое уравнение для функции цены $u(t, x)$:
\begin{equation}
    \frac{\partial u}{\partial t} + \frac{1}{2}\text{Tr}(\sigma\sigma^T \text{Hess}_x u) + \nabla u \cdot \mu + f(t, x, u, \sigma^T \nabla u) = 0,
\end{equation}
с терминальным условием $u(T, x) = g(x)$. Здесь $f(\cdot)$ — нелинейная функция, описывающая, например, разницу процентных ставок или риск контрагента.

Согласно нелинейной версии формулы Фейнмана-Каца (теорема Парду-Пенга), решение этого PDE эквивалентно решению системы стохастических уравнений:

\begin{enumerate}
    \item \textbf{Прямое уравнение (Forward SDE) — Динамика активов:}
    \begin{equation}
        X_t = \xi + \int_{0}^{t} \mu(s, X_s) ds + \int_{0}^{t} \sigma(s, X_s) dW_s
    \end{equation}
    \item \textbf{Обратное уравнение (Backward SDE) — Динамика портфеля:}
    \begin{equation}
        Y_t = g(X_T) - \int_{t}^{T} f(s, X_s, Y_s, Z_s) ds - \int_{t}^{T} Z_s^T dW_s
    \end{equation}
\end{enumerate}

\textbf{Финансовый смысл переменных:}
\begin{itemize}
    \item $Y_t = u(t, X_t)$ — цена дериватива (или стоимость портфеля) в момент $t$.
    \item $Z_t = \sigma^T(t, X_t) \nabla u(t, X_t)$ — вектор хеджирования (аналог «Дельты», умноженной на волатильность), показывающий, сколько риска нужно перекрыть.
\end{itemize}

\subsubsection{Алгоритм и применение нейросетей}
В классических численных методах $Z_t$ неизвестно. Идея Deep BSDE заключается в аппроксимации неизвестной функции градиента $Z_t$ нейронной сетью на каждом временном шаге.

Алгоритм дискретизуется по времени $0 = t_0 < t_1 < \dots < t_N = T$ (схема Эйлера-Маруямы):
\begin{equation}
    Y_{n+1} \approx Y_n - f(t_n, X_n, Y_n, Z_n)\Delta t + Z_n^T \Delta W_n
\end{equation}
Здесь $Z_n = \text{NeuralNet}(X_n; \theta_n)$.
Процесс начинается с $Y_0$ (искомая цена, обучаемый параметр) и идет вперед во времени.

Функция потерь (Loss Function) строится на невязке терминального условия:
\begin{equation}
    L(\theta) = \mathbb{E} \left[ | Y_T(\theta) - g(X_T) |^2 \right]
\end{equation}
Минимизируя разницу между тем, что мы «насимулировали» ($Y_T$), и тем, что должны выплатить по контракту ($g(X_T)$), сеть одновременно находит справедливую цену $Y_0$ и стратегию хеджирования $\{Z_t\}$.

\subsubsection{Область применения и преимущества}
\begin{itemize}
    \item \textbf{Где используется:} Основная ниша метода — задачи оценки кредитных рисков (\textbf{XVA: CVA/DVA/FVA}) для крупных портфелей банков, а также прайсинг опционов на корзины акций (Basket Options) в условиях нелинейной волатильности или ставок.
    \item \textbf{Кем используется:} Исследовательские подразделения (Quant Research) крупных инвестиционных банков (JP Morgan, Goldman Sachs) и хедж-фонды, работающие со сложными структурными продуктами.
    \item \textbf{Когда применять в дипломе:}
    \begin{itemize}
        \item Если размерность задачи $d \ge 10$ (классические сетки умирают).
        \item Если задача нелинейная (обычный Монте-Карло не подходит).
        \item Если нужно найти не только цену, но и хеджирующую стратегию.
    \end{itemize}
    \item \textbf{Ограничения:} В отличие от метода DGM, Deep BSDE находит решение только для \textit{одного} фиксированного начального условия $(t=0, X_0)$. Если цена актива $X_0$ изменится, сеть нужно переобучать заново (или использовать методы параметризации $X_0$, что сложнее).
\end{itemize}

% --- СТАТЬЯ 2: SIRIGNANO & SPILIOPOULOS (DGM) ---
\subsection{Метод DGM: Глобальное решение и задачи со свободной границей}
\textbf{Статья:} \textit{DGM: A deep learning algorithm for solving partial differential equations} \cite{Sirignano2018}.

\subsubsection{Проблема и мотивация}
В то время как метод Deep BSDE решает уравнение вдоль конкретных стохастических траекторий (Lagrangian approach), он имеет существенный недостаток: решение находится только для одной фиксированной точки $(t=0, X_0)$. Если рыночная цена актива изменится, процедуру обучения нужно запускать заново.

Метод DGM (Deep Galerkin Method) предлагает \textit{глобальный} подход (Eulerian approach). Цель метода — обучить нейросеть, которая будет аппроксимировать функцию цены $u(t, x)$ во всей области определения сразу. Это позволяет не только получать мгновенные оценки при изменении цены актива, но и эффективно решать задачи с \textbf{свободной границей} (Free Boundary Problems), к которым относится оценка Американских опционов — наиболее распространенного типа деривативов на биржевых рынках.

\subsubsection{Математическая теория}
Рассмотрим общее параболическое PDE для функции $u(t, x)$ в области $\Omega \subset \mathbb{R}^d$ на интервале $[0, T]$:
\begin{equation}
    (\partial_t + \mathcal{L})u(t, x) = 0, \quad (t, x) \in [0, T] \times \Omega,
\end{equation}
где $\mathcal{L}$ — дифференциальный оператор (для финансов — оператор Блэка-Шоулза или Хестона), включающий вторые производные (Гессиан). Также заданы граничные условия на $\partial \Omega$ и терминальное условие при $t=T$.

Вместо построения сетки, авторы предлагают искать решение в виде нейронной сети $f(t, x; \theta)$. Функция потерь конструируется как взвешенная сумма квадратов невязок (residuals) уравнения в случайно выбранных точках:

\begin{equation}
    J(f) = \underbrace{|| \partial_t f + \mathcal{L}f ||^2_{[0,T]\times \Omega}}_{\text{Ошибка внутри области}} + \underbrace{||f - g||^2_{\partial \Omega \times [0, T]}}_{\text{Граничные условия}} + \underbrace{||f(T, \cdot) - \text{Payoff}||^2_{\Omega}}_{\text{Терминальное условие}}
\end{equation}

\textbf{Ключевая математическая инновация:}
Вычисление оператора $\mathcal{L}f$ требует нахождения вторых производных (матрицы Гессиана). Для нейросети в высокой размерности $d$ это операция сложности $O(d^2)$, что вычислительно запретительно. Авторы предложили метод несмещенной оценки вторых производных с помощью Монте-Карло (аналог Hutchinson's trace estimator), что снижает сложность до $O(d)$ и позволяет работать с размерностями $d=200$.

\subsubsection{Алгоритм и Архитектура}
Процесс обучения (Training Loop) выглядит следующим образом:
\begin{enumerate}
    \item На каждой итерации генерируется пакет (batch) случайных точек внутри области, на границах и во времени $T$.
    \item С помощью автоматического дифференцирования (Autograd) вычисляются производные нейросети по входам $\frac{\partial f}{\partial t}, \frac{\partial f}{\partial x}, \frac{\partial^2 f}{\partial x^2}$.
    \item Вычисляется невязка уравнения и обновляются веса $\theta$ методом стохастического градиентного спуска (Adam).
\end{enumerate}

Архитектура сети, предложенная авторами, вдохновлена LSTM и ResNet (DGM-Layer). Она включает проброс связей (skip-connections) и мультипликативные взаимодействия, что позволяет избежать затухания градиента при вычислении производных высоких порядков.

\subsubsection{Область применения и преимущества}
\begin{itemize}
    \item \textbf{Американские опционы:} Это главная "фишка" метода. Задача прайсинга Американского опциона формулируется как задача с препятствием (obstacle problem) или задача со свободной границей. В DGM это реализуется добавлением штрафа в Loss-функцию, который гарантирует, что цена опциона всегда не ниже выплаты при досрочном исполнении. Deep BSDE справляется с этим значительно хуже (требует Reflected BSDE).
    \item \textbf{Универсальность:} Одна обученная модель дает поверхность цен для любых $S_t$ и $t$.
    \item \textbf{Сложность вычислений:} Метод является вычислительно тяжелым («expensive»). Взятие вторых производных через граф вычислений нейросети требует большого объема видеопамяти (VRAM) и мощных GPU.
    \item \textbf{Когда применять в дипломе:}
    \begin{itemize}
        \item Если тема касается \textbf{Американских опционов} или Бермудских опционов.
        \item Если нужно исследовать зависимость цены от состояния рынка (построить график Value Surface).
        \item Если имеется доступ к хорошим вычислительным ресурсам (NVIDIA Tesla T4/P100 и выше).
    \end{itemize}
\end{itemize}

% =================================================================
% РАЗДЕЛ 3: ДЕТАЛЬНЫЙ АНАЛИЗ КЛЮЧЕВЫХ ИССЛЕДОВАНИЙ
% =================================================================
\section{Детальный анализ ключевых исследований}

В данном разделе проводится анализ фундаментальных работ, формирующих методологическую базу применения глубокого обучения в задачах финансовой математики. Рассматриваются методы решения уравнений в частных производных (PDE), подходы к хеджированию в условиях рыночных несовершенств, а также алгоритмы калибровки моделей.

% --- 3.1 DEEP BSDE ---
\subsection{Метод Deep BSDE: Стохастическая оптимизация для высокоразмерных PDE}
\textbf{Статья:} \textit{Solving high-dimensional partial differential equations using deep learning} \cite{Han2018}.

\subsubsection{Проблема и мотивация}
Классические сеточные методы (Finite Difference Methods) подвержены «проклятию размерности»: количество узлов сетки растет экспоненциально с увеличением числа активов, что делает их неприменимыми для размерностей $d > 5$. Стандартный метод Монте-Карло, работающий в высоких размерностях, применим преимущественно для линейных PDE.
Однако многие современные финансовые задачи (например, оценка CVA, прайсинг с учетом дефолтного риска или различных ставок заимствования) описываются \textbf{нелинейными} PDE. Метод Deep BSDE стал первым алгоритмом, позволившим решать такие уравнения в размерности $d=100$ за полиномиальное время.

\subsubsection{Математическая теория}
Метод базируется на теории обратных стохастических дифференциальных уравнений (Backward Stochastic Differential Equations, BSDE). Рассмотрим полулинейное параболическое уравнение для функции цены $u(t, x)$:
\begin{equation}
    \frac{\partial u}{\partial t} + \frac{1}{2}\text{Tr}(\sigma\sigma^T \text{Hess}_x u) + \nabla u \cdot \mu + f(t, x, u, \sigma^T \nabla u) = 0,
\end{equation}
с терминальным условием $u(T, x) = g(x)$.

Согласно нелинейной версии формулы Фейнмана-Каца, решение этого PDE эквивалентно решению системы стохастических уравнений:
\begin{equation}
    Y_t = g(X_T) - \int_{t}^{T} f(s, X_s, Y_s, Z_s) ds - \int_{t}^{T} Z_s^T dW_s,
\end{equation}
где $Y_t = u(t, X_t)$ — цена дериватива, а $Z_t = \sigma^T \nabla u(t, X_t)$ — вектор хеджирования.

\subsubsection{Алгоритм и применимость}
Идея метода заключается в аппроксимации неизвестной функции градиента $Z_t$ нейронной сетью на каждом временном шаге дискретизации Эйлера-Маруямы. Функция потерь строится на невязке терминального условия:
\begin{equation}
    L(\theta) = \mathbb{E} \left[ | Y_T(\theta) - g(X_T) |^2 \right].
\end{equation}
Данный подход является стандартом (SOTA) для задач высокой размерности (например, оценка корзинных опционов). Его применение целесообразно в случаях, когда требуется найти не только цену, но и стратегию хеджирования в условиях нелинейной динамики.

% --- 3.2 DGM ---
\subsection{DGM: Бессеточный метод Галёркина}
\textbf{Статья:} \textit{DGM: A deep learning algorithm for solving partial differential equations} \cite{Sirignano2018}.

\subsubsection{Проблема и мотивация}
Метод Deep BSDE решает уравнение вдоль конкретных стохастических траекторий, что ограничивает решение одной фиксированной точкой $(t=0, X_0)$. Метод DGM (Deep Galerkin Method) предлагает глобальный подход, обучая нейросеть аппроксимировать функцию цены во всей области определения. Это позволяет эффективно решать задачи со свободной границей (Free Boundary Problems), к которым относится оценка Американских опционов.

\subsubsection{Математическая теория}
Обучение сводится к минимизации квадрата невязки дифференциального оператора $\mathcal{L}$ в случайно выбранных точках области $\Omega$:
\begin{equation}
    J(f) = || \partial_t f + \mathcal{L}f ||^2_{\Omega} + ||f - g||^2_{\partial \Omega} + ||f(T, \cdot) - \text{Payoff}||^2.
\end{equation}
Авторы предложили метод несмещенной оценки вторых производных (Гессиана) с помощью Монте-Карло, что снижает вычислительную сложность с $O(d^2)$ до $O(d)$.

\subsubsection{Алгоритм и применимость}
Архитектура сети, основанная на DGM-слоях (модификация ResNet/LSTM), позволяет избежать затухания градиента при вычислении производных высоких порядков. Метод наиболее эффективен для задач, требующих построения полной поверхности цен (Value Surface) и оценки опционов с возможностью досрочного исполнения, однако требует значительных вычислительных ресурсов (GPU) для обучения.

% --- 3.3 DEEP HEDGING ---
\subsection{Deep Hedging: Оптимизация торговых стратегий при наличии трений}
\textbf{Статья:} \textit{Deep Hedging} \cite{Buehler2019}.

\subsubsection{Проблема и мотивация}
Классическая теория (Блэк-Шоулз) предполагает отсутствие транзакционных издержек и возможность непрерывной торговли. В реальных условиях эти допущения нарушаются, делая классическое дельта-хеджирование субоптимальным. Подход \textit{Deep Hedging} предлагает смену парадигмы: от поиска «теоретической цены» к поиску оптимальной стратегии, минимизирующей риск с учетом рыночных трений.

\subsubsection{Математическая теория}
Задача формулируется как проблема стохастического управления. Итоговый капитал $PL_T$ зависит от стратегии $\delta$ и транзакционных издержек $C_k$, которые зависят от оборота $|\delta_k - \delta_{k-1}|$:
\begin{equation}
    PL_T(\delta) = -Z + p_0 + \sum (\delta \cdot \Delta S) - \sum C_k(\delta_k, \delta_{k-1}).
\end{equation}
Целью является минимизация выпуклой меры риска $\rho$ (например, CVaR или экспоненциальной полезности):
\begin{equation}
    \pi(Z) = \inf_{\delta \in \mathcal{H}} \rho \left( - PL_T(\delta) \right).
\end{equation}

\subsubsection{Алгоритм и применимость}
Поскольку издержки зависят от предыдущего состояния портфеля, задача является немарковской относительно цены актива. Для решения используются рекуррентные нейронные сети (RNN/LSTM). Метод демонстрирует высокую практическую ценность, позволяя моделировать стратегии, превосходящие классические подходы по показателю риск/доходность на рынках с комиссиями.

% --- 3.4 DEEP CALIBRATION ---
\subsection{Deep Calibration: Ускорение калибровки сложных моделей}
\textbf{Статья:} \textit{Deep learning volatility: a deep neural network perspective on pricing and calibration in (rough) volatility models} \cite{Horvath2019}.

\subsubsection{Проблема и мотивация}
Современные модели, такие как \textit{Rough Volatility}, точно описывают динамику рынка, но не имеют аналитических решений. Их калибровка традиционными методами требует множественных симуляций Монте-Карло, что занимает неприемлемо много времени для операционной деятельности.

\subsubsection{Суть метода}
Предлагается двухэтапный подход:
\begin{enumerate}
    \item \textbf{Офлайн-обучение:} Нейросеть обучается на синтетических данных аппроксимировать функцию цены $F(\theta) \approx P(\theta)$. Сеть выступает в роли детерминированного суррогата стохастической модели.
    \item \textbf{Онлайн-калибровка:} При поступлении рыночных данных параметры модели $\theta^*$ находятся путем минимизации разницы между предсказанием сети и рынком с использованием градиентного оптимизатора.
\end{enumerate}

\subsubsection{Алгоритм и применимость}
Использование суррогатной модели позволяет ускорить процесс калибровки в $10\,000$--$60\,000$ раз. Метод применим к широкому классу моделей (Хестон, Rough Bergomi) и является эффективным решением проблемы «вычислительного бутылочного горлышка» в финансовом моделировании.

% --- 3.5 CaNN ---
\subsection{CaNN: Калибровка через механизм обратного распространения ошибки}
\textbf{Статья:} \textit{A neural network-based framework for financial model calibration} \cite{Liu2019}.

\subsubsection{Суть метода}
В отличие от подхода с внешним оптимизатором, фреймворк CaNN (Calibration Neural Network) использует механизм обратного распространения ошибки (Backpropagation) для решения обратной задачи.
Процесс разделен на две фазы: обучение сети (forward pass) и калибровка (backward pass). На этапе калибровки веса сети замораживаются, а входной слой, соответствующий параметрам модели, объявляется обучаемым. Градиент ошибки между предсказанной и рыночной ценой распространяется до входа, обновляя параметры модели.

\subsubsection{Применимость}
Данный подход позволяет выполнять калибровку, используя встроенные возможности фреймворков глубокого обучения (например, PyTorch) без подключения сторонних солверов, что упрощает архитектуру программного решения.

% --- 3.6 ТЕОРЕТИЧЕСКОЕ ОБОСНОВАНИЕ ---
\subsection{Теоретическое обоснование: Рандомизированные сигнатуры}
\textbf{Статья:} \textit{Discrete-time signatures and randomness in reservoir computing} \cite{Cuchiero2021}.

\subsubsection{Суть метода}
Статья предоставляет фундаментальное теоретическое обоснование эффективности рекуррентных нейронных сетей (RNN) в задачах прогнозирования финансовых временных рядов. Авторы связывают RNN с теорией сигнатур путей (Path Signatures) и рядами Вольтерра, доказывая свойство универсальной аппроксимации для фильтров с затухающей памятью.

\subsubsection{Применимость}
Работа обосновывает выбор архитектур типа LSTM/GRU для задач, зависящих от траектории (path-dependent), и предлагает альтернативные методы на основе Reservoir Computing, которые могут быть вычислительно эффективнее полного обучения глубоких сетей.

% --- 3.7 ИНЖЕНЕРНЫЕ АСПЕКТЫ ---
\subsection{Инженерные аспекты: Обеспечение финансовой корректности}
\textbf{Статья:} \textit{Deep learning calibration of option pricing models: some pitfalls and solutions} \cite{Itkin2019}.

\subsubsection{Проблема и решения}
Прямое применение стандартных методов ML может приводить к моделям, допускающим статический арбитраж или имеющим разрывные производные (Греки).
Автор предлагает ряд инженерных решений:
\begin{enumerate}
    \item Использование гладких функций активации класса $C^2$ (Softplus, ELU) вместо ReLU для корректного вычисления вторых производных (Гамма).
    \item Введение штрафов (Soft Constraints) в функцию потерь за нарушение условий отсутствия арбитража (например, нарушение монотонности цены по времени или выпуклости по страйку).
\end{enumerate}

\subsubsection{Применимость}
Данные рекомендации являются критически важными при практической реализации моделей для обеспечения их робастности и соответствия финансовой теории.
% =================================================================
% СРАВНЕНИЕ И ВЫВОДЫ
% =================================================================
\section{Сравнительный анализ и Выводы}

\begin{table}[h!]
    \centering
    \caption{Сравнение нейросетевых подходов}
    \label{tab:comparison}
    \begin{tabular}{|p{3cm}|p{4cm}|p{3cm}|p{4cm}|}
        \hline
        \textbf{Метод} & \textbf{Задача} & \textbf{Сложность} & \textbf{Ресурсы (GPU)} \\
        \hline
        Deep BSDE & Прайсинг, $d \sim 100$ & Средняя & Средние (T4) \\
        \hline
        DGM & Американские опционы & Высокая & Высокие (V100/A100) \\
        \hline
        Deep Hedging & Хеджирование с трениями & Средняя & Средние/Высокие \\
        \hline
        Deep Calibration & Мгновенная калибровка & Низкая & Низкие (любой GPU) \\
        \hline
    \end{tabular}
\end{table}

\subsection{План практического исследования}
На основе проведенного обзора, для практической части дипломной работы выбран метод \textbf{Deep Calibration}. Данный подход позволяет решить актуальную проблему медленной калибровки, демонстрируя превосходство над классическими методами.
Для реализации будет использован фреймворк PyTorch, а обучение будет проводиться с использованием облачных ресурсов (Google Colab/Kaggle) для ускорения вычислений.

\newpage
\printbibliography

\end{document}